\section{The annotation model}
\label{sec:annotations}

\subsection{Geometry}

An annotation stores its geometry in unrotated page space at scale~1, origin top-left, $y$
increasing downwards. Four consequences follow, all of them the reason for the choice:

\begin{itemize}
  \item a zoom changes a scale factor and nothing else;
  \item a rotation changes the transform used to paint, not the coordinates stored;
  \item a view-mode change (single, continuous, facing, facing-continuous) does not touch the
        model at all;
  \item hit testing compares a point in the same space, so it is a rectangle test rather than
        an inverse projection.
\end{itemize}

PDF user space is the mirror image, with an origin at the bottom-left and a scale given by the
page's \texttt{/UserUnit}, so writing an annotation is an explicit
\texttt{toScreen}/\texttt{fromScreen} conversion in one place
(\texttt{src/annotations/layer.ts}). Keeping that conversion in one function, rather than
inline at each use, is what made it possible to add rotated pages without auditing a dozen
call sites.

\subsection{Kinds and the dual representation}
\label{sec:annot-model}

Fifteen kinds are modelled. Twelve of them draw or carry text: highlight, underline,
strikeout, \emph{replace text}, freehand ink, sticky note, rectangle, ellipse, arrow,
typewriter (the original product's name for a free text box), callout text, and stamp --- where
stamp takes one of six standard names. The remaining three are measurements: distance,
perimeter and area.

Each is written as a standard subtype, so that other readers display it:
\texttt{/Highlight}, \texttt{/Underline}, \texttt{/StrikeOut}, \texttt{/Ink}, \texttt{/Text},
\texttt{/Square}, \texttt{/Circle}, \texttt{/Line}, \texttt{/FreeText}, \texttt{/PolyLine},
\texttt{/Polygon}, \texttt{/Stamp}. Alongside, a private \texttt{/PLP} dictionary carries
\texttt{/K} (the exact kind), \texttt{/W} (stroke width), \texttt{/D} (dash) and
\texttt{/F} (font size). Readers that do not know \texttt{/PLP} ignore it; this one prefers it
on load and falls back to inferring a kind from the subtype, which is how annotations made in
other applications are displayed with sensible defaults.

\paragraph{The cost of two representations.} This duality is the model's weakest point and
worth stating plainly. A third-party \texttt{/StrikeOut} with no \texttt{/PLP} becomes a
strikeout in this reader, which is correct; a \texttt{/StrikeOut} written by this reader for a
\emph{replace text} annotation becomes a strikethrough everywhere else, which is the closest
the format can express and is what the paper's \S\ref{sec:pdf-layer} describes as a
limitation rather than a feature. Two representations can also disagree --- a file edited in
another application may carry a \texttt{/PLP} kind that its subtype no longer supports --- and
the loader resolves that in favour of the subtype when the two are incompatible, silently. We
know of no such file; the rule exists because guessing in the other direction would be worse.

\paragraph{Threads.} A reply is its own annotation, a \texttt{/Text} object whose \texttt{/IRT}
points at its parent through an indirect reference; this is the mechanism every reader uses
to display a conversation. Reading back therefore takes two passes: a reply may reference a
parent that appears later in the page's list, so the loader collects replies first and attaches
them once every annotation is known. A reply whose parent lives on another page is still
shown, rather than dropped, because a missing parent is a fact about the file rather than
about the reply.

\paragraph{Identity and deletion.} Annotations carry a locally generated identifier and,
separately, the object number they were read from. The two must not be conflated: the
identifier survives an edit, while the object number is what a save must remove. Keeping the
second only on the document session, not on the annotation, is the design that made
\S\ref{sec:incremental} correct.

\subsection{Interaction}

Gestures are owned by the annotation layer rather than by the tool manager, because they
differ per kind: text markup consumes a selection, ink and rectangles drag out geometry, a
note is a single click, and a stamp is a click with a size of its own.

Two interaction rules were learned from defects:

\begin{itemize}
  \item \textbf{A click selects an annotation under any tool}, and the \emph{gesture} belongs
        to the active tool. An earlier version stopped propagation whenever a pointer landed on
        an annotation, which broke drag-to-scroll and marquee selection for gestures that
        merely started on top of a mark.
  \item \textbf{A text tool pressed on an existing text box edits it} rather than stacking a
        blank box on top. The earlier behaviour selected the annotation and then covered it with
        an empty box, which a user reasonably reported as ``I cannot select my annotation''.
\end{itemize}

\subsection{Rendering and printing}

Pages are virtualised: only pages near the viewport have a rasterised canvas, and a page
leaving the window drops its canvas while keeping its element. \S\ref{sec:eval} shows the
consequence in the measurements --- opening a 200-page document and a 24-page document take
the same time to first painted page.

Printing rasterises the selected range at twice the page size (144\,DPI) into a hidden
container, hides the application shell with a print stylesheet, and calls the platform's print
dialog; in the desktop build the same is done through the Electron main process. The printed
sheets are removed when the print job reports completion rather than on a timer, because a
timer loses the race on a slow printer.
